% =====================================================================
% Analisi Matematica II -- Lezione 4 (aa)
% Derivazione per serie (le serie di potenze sono in serie-di-potenze.tex).
% Fonti: appunti manoscritti (pp. 1--2) + audio "Analisi II - Lezione 4".
% =====================================================================

% [0:00:32]
\section{Teorema di derivazione per serie}

% [0:00:00]
Passando dalle successioni di funzioni (convergenza puntuale e uniforme) alle
serie di funzioni i tipi di convergenza sono diventati quattro: puntuale,
uniforme, assoluta e totale. Come per le successioni valgono il teorema di
continuità del limite e quello di passaggio al limite sotto il segno di
integrale, per le serie valgono la continuità della somma e l'integrazione per
serie. % [0:00:32]
Manca l'analogo del passaggio al limite sotto il segno di derivata: la
\emph{derivazione per serie}.

\paragraph{Teorema (derivazione per serie).} Siano $f_n\colon[a,b]\to\mathbb{R}$,
$n\in\mathbb{N}$, tali che:
\begin{enumerate}
  \item $f_n\in C^1([a,b])$ per ogni $n$;
  \item la serie converge puntualmente in $[a,b]$:
  \[
    \sum_{n=0}^{\infty} f_n(x)\xrightarrow{\;P\;} S(x)\quad\text{in }[a,b];
  \]
  \item la serie delle derivate converge uniformemente in $[a,b]$:
  \[
    \sum_{n=0}^{\infty} f_n'(x)\xrightarrow{\;U\;} G(x)\quad\text{in }[a,b].
  \]
\end{enumerate}
Allora
\[
  S\in C^1([a,b]),\qquad S'(x)=G(x)\quad\forall x\in[a,b],
\]
cioè si può derivare termine a termine:
\[
  \frac{d}{dx}\sum_{n=0}^{\infty} f_n(x)=\sum_{n=0}^{\infty} f_n'(x).
\]
Inoltre (è implicito) la successione delle somme parziali converge
uniformemente:
\[
  S_n(x)=\sum_{k=0}^{n} f_k(x)\xrightarrow{\;U\;} S(x)\quad\text{in }[a,b].
\]

\paragraph{Dimostrazione.} % [0:05:16]
È quella del teorema di passaggio al limite sotto il segno di derivata per le
successioni di funzioni, applicato alla successione delle somme parziali $S_n$.

\paragraph{Osservazione (perché questo enunciato).}
% [0:01:51]
Il teorema per le successioni chiedeva la convergenza in un solo punto $x_0$;
qui si chiede la convergenza puntuale in tutto $[a,b]$. L'enunciato sembra più
debole, ma è scritto nella forma in cui servirà per le serie di potenze.
% [0:05:16]
Il motivo è pratico: per una successione la terza ipotesi si verifica a mano,
mentre per una serie la convergenza uniforme delle derivate si ottiene in
pratica solo tramite la convergenza totale, cioè cercando una $M_n$
indipendente da $x$. Così si studia la serie in tutte le $x$, e in genere anche
la serie di partenza converge già ovunque: tanto vale chiederne la convergenza
puntuale in tutto $[a,b]$.
% [0:04:02]
La convergenza uniforme delle $S_n$ nella tesi non è un fatto nuovo: è quella
fornita dal teorema sulle successioni.

